\section{Supplemental requirements}
In addition to the features discovered in the use cases, we added some
requirements we felt were needed, that can be broken down into a RUP
FURPS+ fashion \cite{FURPS}.
\subsection{Functionality}
	\begin{itemize}
		\item The system should be able to restart a \snort{} sensor in the event that
		one crashes.
		\item The system should be able to detect if a \snort{} sensor is not responsive
		over the network.
		\item The system must be portable to the majority of the popular Linux
		distributions. Porting to other operating systems, such as Windows, should be
		possible through changing the operating system-specific code.
	\end{itemize}
\subsection{Usability}
	\begin{itemize}
		\item The graphical user interface must be simple and clean, but still contain
		all the features any user might require to interact with the back-end system.
		\item The system interfaces, herein API and function calls, must be
		self-explanatory.
	\end{itemize}
\subsection{Reliability}
	\begin{itemize}
		\item The system must not cause a \snort{} sensor to crash.
		\item The system must be responsive to API-calls at all times. In the event
		the system capacity limit is reached, new requests should be queued.
	\end{itemize}
\subsection{Performance}
	\begin{itemize}
		\item The system should handle administration of at least 200 sensors.
		\item The database should be able to maintain at least 100 000 individual rules.
		\item Any operation in the system in normal operations (i.e. after the first big
		rule update) should take no longer than 1 second on average.
		\item The relationship between time and workload should be sublinear, meaning
		that if a single operation takes 1ms, ten operations should take 10ms or less.
	\end{itemize}
\subsection{Security}
	\begin{itemize}
		\item The link between the system server and the remote sensor must be secured in such a way that a man in the middle cannot read the transmission in plaintext. 
		\item The remote sensor and the system server must verify that they are trusted
		parties in the system, through authentication.
		\item Certain important API-calls must be protected against misuse from a
		non-trusted party.
		\item All processes and files related to the system must exhibit the principle
		of least privilege \cite{saltzer}.
	\end{itemize}
\subsection{Interoperability}
The system must be able to work with \snort{} in such a way that it can
import new files to the \snort{} configuration files structure and must be able to
reset \snort{} in such a way that the new configurations are loaded.
\subsection{Licensing}
The system will be released to the public as Open Source under the GNU General
Public License v3 license\cite{GPL}.
